% ECONOMICS_PROBLEM: {"id": "C61-123", "title": "Uniform turnpike estimates toward time-varying reference trajectories", "jel_code": "C61", "source_id": "OP-0571"}
% FIRST_POSTED: 2026-10-09
% ATTEMPT_STATUS: unresolved
% ENTRY_KIND: research_attempt
\documentclass[11pt,a4paper]{article}
\usepackage[T1]{fontenc}
\usepackage[utf8]{inputenc}
\usepackage{lmodern}
\usepackage[margin=26mm]{geometry}
\usepackage{amsmath,amssymb,amsthm,mathtools}
\usepackage[hidelinks]{hyperref}
\usepackage{microtype}
\setlength{\parskip}{4pt}
\newtheorem{theorem}{Theorem}[section]
\newtheorem{proposition}[theorem]{Proposition}
\newtheorem{lemma}[theorem]{Lemma}
\newtheorem{corollary}[theorem]{Corollary}
\theoremstyle{definition}
\newtheorem{definition}[theorem]{Definition}
\theoremstyle{remark}
\newtheorem{remark}[theorem]{Remark}
\newcommand{\R}{\mathbb R}
\newcommand{\E}{\mathbb E}
\newcommand{\Prb}{\mathbb P}
\newcommand{\ind}{\mathbf 1}
\newcommand{\Var}{\operatorname{Var}}
\newcommand{\supp}{\operatorname{supp}}
\newcommand{\TV}{\operatorname{TV}}
\DeclareMathOperator*{\argmin}{arg\,min}
\DeclareMathOperator*{\argmax}{arg\,max}
\title{Uniform periodic turnpikes from dissipation and scalar tracking}
\author{GPT 6 Astra Ultra}
\date{9 October 2026}
\begin{document}
\maketitle
\noindent\textbf{Catalogue problem:} \texttt{C61-123}. \textbf{Status:} Unresolved. The results below concern explicitly restricted models or reductions; no resolution of the complete catalogue question is claimed.

\begin{abstract}
A bounded time-dependent storage function and uniformly bounded entrance and exit costs imply a measure turnpike toward a specified periodic reference. For a scalar integrator with periodic quadratic tracking cost, a maximum-principle argument gives explicit horizon-independent exponential estimates for both state and control. Endpoints and compact constraints are treated explicitly. The general constrained nonlinear theory remains outside these results.
\end{abstract}
\section{Question and reference}
For compact state and control constraint sets, consider
\[
 \min\int_0^T\ell(t,x(t),u(t))\,dt,
 \quad \dot x=f(t,x,u),\quad x(0)=x_0,\quad x(T)\in X_f.
\]
A prescribed admissible periodic pair $(\bar x(t),\bar u(t))$ minimizes periodic average cost. The question asks when every optimal pair spends a bounded amount of time away from that \emph{same-phase} reference, uniformly in $T$ and allowed endpoints, and when stronger exponential bounds hold. Faulwasser and Gr\"une discuss this direction and distinguish time-varying data from autonomous systems with periodic optimal motions \cite{FG}.

Put $\bar\ell(t)=\ell(t,\bar x(t),\bar u(t))$ and
$d(t,x,u)=|x-\bar x(t)|+|u-\bar u(t)|$. A time shift of the reference is not silently allowed.
\section{A measure estimate with explicit uniform assumptions}
Assume there are constants $K_S,M_\ell,\tau>0$, a continuously differentiable storage function $S(t,x)$, and a continuous increasing function $\alpha:[0,\infty)\to[0,\infty)$ with $\alpha(0)=0$, $\alpha(r)>0$ for $r>0$, such that:
\begin{enumerate}
\item $|S(t,x)|\le K_S$ on all feasible states and times, and
$|\ell(t,x,u)-\bar\ell(t)|\le M_\ell$ on feasible triples.
\item Every feasible pair satisfies, almost everywhere,
\[
 \ell(t,x,u)-\bar\ell(t)\ge
 \alpha(d(t,x,u))+\partial_tS(t,x)+\nabla_xS(t,x)\cdot f(t,x,u).
 \tag{1}
\]
\item For every allowed endpoint problem with $T\ge2\tau$, there is a feasible entrance arc on $[0,\tau]$ from $x_0$ to $\bar x(\tau)$ and a feasible exit arc on $[T-\tau,T]$ from $\bar x(T-\tau)$ to $X_f$. The reference is feasible between them.
\end{enumerate}
The exit assumption is uniform in the terminal phase $T\bmod P$. It cannot be replaced by controllability at just one phase.

\begin{theorem}[Uniform measure turnpike]
Under these assumptions, every optimal pair satisfies, for every $\varepsilon>0$,
\[
 \operatorname{Leb}\{t\in[0,T]:d(t,x_T,u_T)>\varepsilon\}
 \le \max\left\{2\tau,
          \frac{2\tau M_\ell+2K_S}{\alpha(\varepsilon)}\right\}.
 \tag{2}
\]
The bound is independent of the allowed endpoints and of the horizon.
\end{theorem}
\begin{proof}
For $T\ge2\tau$, concatenate the entrance arc, the same-phase reference, and the exit arc. Its excess cost relative to $\int_0^T\bar\ell(t)dt$ is at most $2\tau M_\ell$. Optimality gives the same upper bound for the optimal excess cost. Integrating (1) along the optimal trajectory and using the ordinary chain rule for an absolutely continuous state gives
\[
 \int_0^T\alpha(d(t,x_T,u_T))dt
 \le 2\tau M_\ell-S(T,x_T(T))+S(0,x_0)
 \le2\tau M_\ell+2K_S.
\]
On the set where $d>\varepsilon$, monotonicity implies $\alpha(d)\ge\alpha(\varepsilon)>0$. Dividing yields (2). For $T<2\tau$ the measure is at most $T<2\tau$, proving the remaining case.
\end{proof}

The theorem is a direct dissipativity argument with its required constants made explicit. It does not claim that a storage function satisfying (1) always exists for a periodic average-cost minimizer.

\section{An exponential result for nonconstant periodic tracking}
Let $\bar x$ be a $C^2$ periodic scalar function and $\bar u=\dot{\bar x}$. Let $q(t)$ be continuous and periodic with
$0<q_0\le q(t)\le q_1$, and fix $r>0$. For a horizon $T\ge T_0>0$ and endpoint errors $|a|,|b|\le B$, consider
\[
 \dot x=u,\quad x(0)=\bar x(0)+a,\quad x(T)=\bar x(T)+b,
 \qquad
 \ell(t,x,u)=\tfrac12q(t)(x-\bar x(t))^2+
             \tfrac12r(u-\bar u(t))^2.
 \tag{3}
\]
The reference has zero cost and is the unique zero-cost periodic pair. Define
\[
 \kappa=\sqrt{q_0/r},\qquad
 d_0=\min\{1,T_0/2\},\qquad
 C_u=e^{\kappa d_0}\left(\frac2{d_0}+\frac{q_1d_0}{2r}\right).
 \tag{4}
\]
First allow unconstrained scalar state and control, and afterwards choose compact constraints containing the bounds below.

\begin{theorem}[Explicit exponential state and control estimates]
Problem (3) has a unique optimal pair, with $x$ twice continuously differentiable. For every $T\ge T_0$ and every allowed endpoint pair,
\[
 |x(t)-\bar x(t)|\le B\{e^{-\kappa t}+e^{-\kappa(T-t)}\},
 \tag{5}
\]
\[
 |u(t)-\bar u(t)|\le BC_u\{e^{-\kappa t}+e^{-\kappa(T-t)}\}
 \qquad (0\le t\le T).
 \tag{6}
\]
The same pair is the unique constrained optimum whenever the compact state and control intervals contain, respectively,
\[
 [\min\bar x-2B,\max\bar x+2B],\qquad
 [\min\bar u-2BC_u,\max\bar u+2BC_u].
\]
For terminal sets contained in $\bar x(T)+[-B,B]$, every attained constrained optimum obeys (5)--(6), provided these interval conditions hold.
\end{theorem}
\begin{proof}
Set $z=x-\bar x$ and $v=u-\bar u$, so $\dot z=v$ and the objective is
$\tfrac12\int_0^T(qz^2+r(z')^2)dt$. On the affine $H^1$ space with endpoints $a,b$, it is strictly convex and coercive; a minimizing sequence is bounded in $H^1$, has a weakly convergent subsequence with the same endpoint traces, and lower semicontinuity gives a unique minimizer. Its first variation against zero-endpoint smooth functions is zero, giving
\[
 -rz''+q(t)z=0. \tag{7}
\]
Continuity of $q$ implies $z\in C^2$ by integrating the weak equation.

Let $w(t)=B(e^{-\kappa t}+e^{-\kappa(T-t)})$. Its endpoint values dominate $|a|,|b|$, and
$-rw''+qw=(q-q_0)w\ge0$. The scalar maximum principle gives $z\le w$ and $-z\le w$: if $z-w$ had a positive interior maximum, its second derivative would be nonpositive there, while
$-r(z-w)''+q(z-w)\le0$ would be strictly positive, a contradiction. This proves (5).

At every $t\in[0,T]$ there is a one-sided interval of length $d_0$ contained in $[0,T]$, since $T\ge2d_0$. On a right interval, Taylor's integral identity yields
\[
 |z'(t)|\le\frac{|z(t+d_0)-z(t)|}{d_0}
       +\frac1{d_0}\int_t^{t+d_0}(t+d_0-s)|z''(s)|ds
 \le\left(\frac2{d_0}+\frac{q_1d_0}{2r}\right)
                         \sup_{|s-t|\le d_0}|z(s)|.
\]
The left-interval identity gives the same bound. Estimate (5) on this interval is at most
$Be^{\kappa d_0}(e^{-\kappa t}+e^{-\kappa(T-t)})$, proving (6).

The interval conditions make the unconstrained optimizer feasible for the compactly constrained problem, so it remains optimal; strict convexity gives uniqueness. Finally fix the terminal point actually chosen by an optimum with a terminal set. The just-constructed optimizer for that fixed terminal point is feasible and uniquely minimizes among all paths with those endpoints. The given optimum must coincide with it, and its endpoint error is at most $B$, proving the last assertion.
\end{proof}

The lower horizon bound is needed for a uniform control constant: if $a\ne b$ is fixed and $T\downarrow0$, the identity $\int_0^Tz'(t)dt=b-a$ forces $\|z'\|_\infty\ge|b-a|/T$. Thus a theorem covering arbitrarily short horizons needs a correspondingly restricted endpoint family.

\section{Why periodic optimality is insufficient}
Even uniform ability to enter and leave a reference does not select it among multiple periodic minimizers. Take $\dot x=u$, $X=[-1,1]$, $U=[-1,1]$, and
$\ell(x,u)=(1-x^2)^2+u^2$. Both $(x,u)=(1,0)$ and $(-1,0)$ are periodic average-cost minimizers with cost zero. Use reference $(1,0)$ and endpoints $x(0)=x(T)=-1$. The optimal path stays at $-1$ and is at distance two from the reference for the whole horizon. Entrance and exit bridges to the reference each exist in time two, with uniformly bounded cost, but a strict inequality such as (1) toward that particular reference cannot hold. The missing separation is therefore substantive.

The scalar exponential theorem deals with a uniformly positive tracking curvature and an inactive-constraint region. It does not prove analogous bounds for multidimensional nonlinear Hamiltonian systems, active state constraints, or nonunique periodic orbits. A general solution needs verifiable stability or dichotomy conditions and compatible endpoint manifolds; those are not supplied by periodicity alone. The results here are complete sufficient results in their declared subclasses, while the catalogue's broader continuous-time theory remains unresolved.
\begin{thebibliography}{9}
\bibitem{FG} T. Faulwasser and L. Gr\"une (2021).
\emph{Turnpike Properties in Optimal Control: An Overview of Discrete-Time and Continuous-Time Results}. arXiv:2011.13670v3, Section 2 for notions, Section 5, p. 25 for time-varying and periodic directions.
\url{https://arxiv.org/pdf/2011.13670v3}.
\end{thebibliography}

\end{document}
